\documentclass[10pt,a4paper]{amsart}
\usepackage[margin=19mm]{geometry}
\usepackage{amsmath,amssymb,mathtools}
\usepackage[scr=rsfs,cal=euler]{mathalfa}
\usepackage{xfrac}
\usepackage[unicode,hidelinks,bookmarks=false]{hyperref}
\newcommand{\ie}{i.e.\ }
\newcommand{\cf}{cf.\ }
\newcommand{\resp}{respectively\ }
\newcommand{\Subsection}{\subsection}
\providecommand{\nobreakdash}{\nobreak\hbox{-}}
\setlength{\emergencystretch}{2em}
\multlinegap0pt
\usepackage{bm}
\usepackage{amssymb}
\usepackage{tikz}
\usepackage{enumerate}
\usepackage[shortlabels]{enumitem}
\usepackage{csquotes}
\usepackage{float}
\usepackage{xcolor}
\usepackage{mathtools}
\newtheorem{cor}{Corollary}
\newtheorem{lemma}{Lemma}
\newcommand{\diam}{\mathrm{diam}}
\newcommand{\qbinom}[2]{\genfrac{(}{)}{0pt}{}{#1}{#2}_q}
\title{Kleitman's theorem over vector spaces: parity phenomena in canonical and global stability}
\author{Chenhui Lv, Zixiang Xu}
\date{Source: arXiv:2605.19526v1 (CC BY 4.0)}
\begin{document}
\maketitle

\noindent\emph{Source and licence.} Kleitman's theorem over vector spaces: parity phenomena in canonical and global stability. Version arXiv:2605.19526v1 (CC BY 4.0), distributed by arXiv under the Creative Commons Attribution 4.0 International licence, http://creativecommons.org/licenses/by/4.0/, as recorded in the arXiv licence metadata for this exact version and shown on the abstract page https://arxiv.org/abs/2605.19526v1. Full licence notice: \texttt{licenses/arxiv-2605.19526-CC-BY-4.0.txt}.\par\smallskip

\noindent\textit{Editorial scope.} Counted unit: the article's lemma lem:equality (source0.tex lines 741--768). The article's complete original proof of it is delivered with it (source0.tex lines 770--832, one \texttt{proof} environment of 63 lines). No mathematics is written, repaired or re-derived: the statement and the proof are the article's own lines, in the article's own notation, and no retained line is substituted or re-flowed.  Retained uncounted as the source-local context the proof needs: lines 707--727.  Nothing was omitted from any retained span, so the package carries no omission record.  No retained line contains a citation, so the wrapper carries no bibliography and no bibliography entry.  No retained line contains a figure environment, an image-inclusion command or drawing code, so no image is delivered and the package declares no asset dependency.  Editorial formatting only, in addition: an AMS \texttt{amsart} preamble that restates the article's own declarations and macros used by the retained lines, namely the environments \texttt{cor}, \texttt{lemma} on separate counters, together with the standard packages those lines need.  The retained environments print in this document as Corollary 1, Lemma 1 in order of appearance, whereas the article prints the same items with the numbering of its own full document.  The complete native source of the article as distributed in the arXiv e-print for this exact version is bundled unchanged as \texttt{sources/200921/source0.tex}.
\begin{cor}\label{lem:complementsum-1}
Let $n\ge d+1$. Let $\mathcal{F} \subseteq \mathcal{V}$ with $\diam(\mathcal{F}) \le d$. For every $0\le k<n-k$, we have
\[
|\mathcal{F}(k)|+|\mathcal{F}(n-k)|
\le
\qbinom{n}{k}.
\]
Moreover, equality holds if and only if either
\(
|\mathcal{F}(k)|=\qbinom{n}{k}\)
and
\(
\mathcal{F}(n-k)=\varnothing,
\)
or
\(
\mathcal{F}(k)=\varnothing\)
and \(
|\mathcal{F}(n-k)|=\qbinom{n}{k}.
\)
\end{cor}

\begin{lemma}\label{lem:equality}
Let $\mathcal{F} \subseteq \mathcal{V}$ with $\diam(\mathcal{F}) \le d$, and put \(t:=\lfloor d/2\rfloor\).
If $n\ge d+1$ and
\(
|\mathcal{F}(k)|+|\mathcal{F}(n-k)|=\qbinom{n}{k}\)
for \(k=0,1,\ldots,t,\)
then either $\mathcal{F}$ or $\mathcal{F}^\perp$ satisfies one of the following:
\begin{enumerate}
\item[\textup{(1)}]
$n=d+1$, and for each $k=0,1,\ldots,t$, either
\(
|\mathcal{F}(k)|=\qbinom{n}{k}\)
and \(
|\mathcal{F}(n-k)|=0,
\)
or vice versa;

\item[\textup{(2)}]
$n\ge d+2$, and
\(
|\mathcal{F}(k)|=\qbinom{n}{k}\)
and
\(
|\mathcal{F}(n-k)|=0
\)
for \(k=0,1,\ldots,t.\)
\end{enumerate}
\end{lemma}

\begin{proof}[Proof of Lemma~\ref{lem:equality}]
Clearly, exactly one of $\mathcal{F}(0)$ and $\mathcal{F}(n)$ is nonempty. 
If $\mathcal{F}(n)\neq \varnothing$, then we pass to $\mathcal{F}^\perp$. 
Thus we may assume that $\mathcal{F}(0)\neq \varnothing$.

By Corollary~\ref{lem:complementsum-1}, for each $0\le k\le t$, either
\(
|\mathcal{F}(k)|=\qbinom{n}{k}\)
and \(
\mathcal{F}(n-k)=\varnothing,
\)
or
\(
\mathcal{F}(k)=\varnothing\) and
\(
|\mathcal{F}(n-k)|=\qbinom{n}{k}.
\)

For part~(1), it suffices to observe that for every $X\in \mathcal{V}(k)$ and every $Y\in \mathcal{V}\setminus \mathcal{V}(n-k)$, we have $\Delta(X,Y)\le d$. 
Indeed, if $\dim Y\le n-k-1$, then
\[
\Delta(X,Y)\le \dim X+\dim Y\le k+(n-k-1)=d.
\]
If $\dim Y\ge n-k+1$, then
\[
\dim(X\cap Y)
=
\dim X+\dim Y-\dim(X+Y)
\ge
k+(n-k+1)-n
\ge 1,
\]
and hence
\[
\Delta(X,Y)
\le
\dim(X+Y)-\dim(X\cap Y)
\le n-1=d.
\]

For part~(2), we argue by induction on $t$. 
The case $t=0$ is immediate. 
Assume the statement holds for $t-1$. Then
\(
|\mathcal{F}(t-1)|=\qbinom{n}{t-1}\)
and
\(
|\mathcal{F}(n-t+1)|=0.
\)
If $\mathcal{F}(n-t)\neq \varnothing$, choose $Y\in \mathcal{F}(n-t)$. Since
\(
|\mathcal{F}(t-1)|=\qbinom{n}{t-1}=|\mathcal{V}(t-1)|,
\)
we may choose $X\in \mathcal{F}(t-1)$ with $X\cap Y= \boldsymbol 0$. Then
\[
\Delta(X,Y)=\dim X+\dim Y=n-1\ge d+1,
\]
since $n\ge d+2$, a contradiction. Therefore $\mathcal{F}(n-t)=\varnothing$, and so
\(
|\mathcal{F}(t)|=\qbinom{n}{t}.
\)
\qedhere
\end{proof}

\end{document}
